\documentclass[11pt]{article}

\usepackage[margin=1.15in]{geometry}
\usepackage{amsmath,amssymb,amsthm}
\usepackage{fontspec}
\usepackage{unicode-math}
\setmathfont{latinmodern-math.otf}
\setmonofont{DejaVu Sans Mono}[Scale=0.82]
\usepackage{xcolor}
\usepackage{booktabs}
\usepackage[colorlinks=true,linkcolor=blue!50!black,citecolor=blue!50!black,
            urlcolor=blue!50!black]{hyperref}

\theoremstyle{plain}
\newtheorem{theorem}{Theorem}
\newtheorem{lemma}[theorem]{Lemma}
\newtheorem{proposition}[theorem]{Proposition}
\newtheorem{corollary}[theorem]{Corollary}
\theoremstyle{definition}
\newtheorem{definition}[theorem]{Definition}
\newtheorem{remark}[theorem]{Remark}

\newcommand{\Prp}{\mathrm{Prop}}
\newcommand{\ess}{\mathrel{\mathsf{ess}}}
\newcommand{\esss}{\mathrel{\mathsf{ess}_{\mathrm{S}}}}
\newcommand{\NE}{\mathsf{NE}}
\newcommand{\NES}{\mathsf{NE}_{\mathrm{S}}}
\newcommand{\Gd}{\mathsf{G}}
\newcommand{\gd}{\mathsf{g}}
\newcommand{\Pos}{\mathsf{P}}

\title{G\"odel's Ontological Argument in Lean 4:\\
       the 1970 inconsistency, modal collapse, and the emendations}
\author{\'Alvaro Lozano-Robledo}
\date{19 September 2026}

\begin{document}
\maketitle

\begin{abstract}
We give a Lean 4 formalization of G\"odel's ontological argument and of the
three emendations of it, with a machine-checked record of exactly what each
result depends on.

G\"odel's own 1970 axioms are \emph{inconsistent}. This was discovered in 2013
by the theorem prover Leo-II and analysed by Benzm\"uller and Woltzenlogel
Paleo~\cite{bwp2016}; the mathematics of Section~\ref{sec:incons} is theirs. The
culprit is a single missing conjunct: G\"odel's definition of \emph{essence}
does not require that $x$ actually have $\varphi$, so the empty property is an
essence of every individual, and Axiom~5 then demands that it be necessarily
instantiated. The refutation needs Axioms 1, 2 and 5 only, and \emph{no frame
condition whatsoever}: it lives in $K$.

Scott's 1987 variant restores the conjunct. We formalize the argument G\"odel
intended --- Theorems 2 and 3, concluding $\Box\exists x\,\Gd(x)$ --- and find
that it needs only \emph{symmetry} of the accessibility relation, the modal
logic $B$, rather than the $S5$ in which it is usually presented. We then
formalize the price: \emph{modal collapse}, $p \to \Box p$, so that nothing is
contingent. Read as a condition on the frame rather than on sentences, collapse
forces the accessibility relation to be \emph{equality} --- every world sees
itself and nothing else --- so $\Box$ and $\Diamond$ become the identity
operation and no modal content survives. That settles the corresponding
question about models: Scott's axioms are satisfiable over a world set of any
size whatsoever, but only as a disjoint heap of isolated one-world models.
There is no non-degenerate model to be found.

All of this presupposes symmetry, and we show that it cannot simply be traded
for the other $S5$ conditions. An
explicit model over a reflexive transitive frame --- an $S4$ frame --- satisfies
Scott's axioms while refuting Theorem 3, modal collapse and the equality of the
frame. So Scott's axioms do not entail modal collapse by themselves; they
entail it in $B$ and $S5$, which is where the argument is always presented. The
price is better read as a price of the frame strength the argument needs than
of the axioms alone.

We then formalize the three emendations that repair this. Anderson's drops the
implausible half of Axiom 1 and strengthens God-likeness and essence to
biconditionals. This turns out to be enough on its own: because God-likeness is
defined by a biconditional whose right-hand side is already a $\Box$, the axiom
asserting God-likeness positive converts in a single step into the
\emph{necessity} of God-likeness, and the essence apparatus is bypassed
entirely. His Theorem 3 then follows in $KB$ from three axioms, his Axiom 4 is
derivable in $K4B$, and his Axiom 5 is not merely derivable but vacuous ---
Anderson's necessary existence holds of every individual at every world. These
last facts are not new; they are reported from automated analysis
in~\cite{bww2017}, and what we add is only that they are now enforced by the
types of the corresponding Lean theorems. Uniqueness, which the biconditional
makes especially direct, needs no frame condition at all --- though it is
available in Scott's system too, from A1b alone. Fitting's makes positivity a
predicate on \emph{extensions} rather than intensional properties; its Theorem 3
needs no frame condition at all, and the reason connects directly to the
previous paragraph. H\'ajek's merges two axioms into one weaker one, and the
resulting definition of God-likeness makes his Axioms 4 and 5 superfluous:
Theorem 3 follows from two of his four axioms. All three avoid modal collapse,
and we exhibit two-world $S5$ models certifying it.

For Fitting's variant we also separate the de re and de dicto readings of the
conclusion. Both are theorems on the full axiom set; on his Part I axioms,
before A4, they come apart, and we give the two-world model that separates
them.

Finally, one small structure does double duty. A model with one world and
one individual satisfies Axioms 1--4 but refutes G\"odel's Axiom 5, certifying
that the inconsistency is not an artefact of our encoding; and it satisfies
\emph{all} of Scott's axioms, certifying that his system is consistent. The
contrast between those two facts --- same model, same notion of positivity,
opposite verdicts on Axiom~5 --- is the sharpest statement of what the missing
conjunct costs.
\end{abstract}

\section{Introduction}

G\"odel's ontological argument is a derivation, in second-order modal logic,
of the conclusion that necessarily there exists a God-like being. It descends
from Leibniz and ultimately from Anselm, but G\"odel's version is unusually
sharp: a primitive predicate $\Pos$ (``is a positive property''), three
definitions, five axioms, three theorems. G\"odel discussed the manuscript with
Dana Scott, who circulated a slightly different version; it is Scott's variant
that most of the subsequent literature analyses, and the two differ in exactly
one conjunct in one definition.

The difference is not cosmetic. In 2013 Benzm\"uller and Woltzenlogel Paleo ran
the routine consistency check that precedes any serious attempt at machine
proof, and the higher-order prover Leo-II reported that G\"odel's axioms were
inconsistent~\cite{bwp2016}. Everything follows from them, including
$\Box \exists x\, \Gd(x)$; the 1970 argument therefore establishes nothing
whatever. G\"odel's omission had been remarked on before --- Hazen called it an
oversight --- but its consequence had gone unnoticed for four decades.

Scott's version is another matter. It is consistent, the argument goes through,
and the objection that has dogged it since Sobel is not unsoundness but
\emph{modal collapse}: the axioms entail that every truth is a necessary truth.

This note formalizes all five systems in Lean 4, in 378 lines of code with no
library dependencies. The point of doing so is not doubt about the mathematics, which
is short and entirely comprehensible, but bookkeeping: in a formalization every
axiom and every frame condition a result consumes appears as a named hypothesis
of that result, and Lean's \verb|#print axioms| makes the foundational
dependencies visible too. Several statements below consequently carry their
hypotheses more precisely than the prose versions do.

\paragraph{What is and is not new here.}
None of the mathematics. The inconsistency, the Empty Essence Lemma and the
refutation of Section~\ref{sec:incons} are due to Benzm\"uller and Woltzenlogel
Paleo~\cite{bwp2016}, whose \S4.1 the derivation follows; Theorems 1--3 are
G\"odel's; modal collapse is Sobel's~\cite{sobel}; the emendations are
Anderson's~\cite{anderson}, Fitting's~\cite{fitting} and
H\'ajek's~\cite{hajek}, and all three were mechanised in Isabelle/HOL
by~\cite{afp,bf2020,bww2017}, from which the definitions used here are taken.
The frame reading of modal collapse (Theorem~\ref{thm:frame}) is standard, and
the superfluousness of H\'ajek's A4 and A5 (Theorem~\ref{thm:hajT3}) is
reported in~\cite{bww2017}.

What this note offers is an independent mechanization in a different
foundation, with a different discipline: axioms and frame conditions are
hypotheses of the theorems that use them rather than global postulates, so the
accounting in Section~\ref{sec:frames} is carried in the statements and
checked; and the models are explicit, constructive proof terms rather than the
output of a model finder. Its purpose is corroboration and exposition, not
discovery --- unlike the work it follows, nothing here was found
automatically. Readers wanting the research frontier should go to the
literature cited above, and to the ongoing dispute over the ultrafilter
reading~\cite{og2026,benz2026}.

The small accounting results of Sections~\ref{sec:frames}, \ref{sec:symm}
and~\ref{sec:derededicto} --- that the 1970 refutation needs no frame
condition, that symmetry is not weakenable towards $S4$, that Scott's axioms
admit models at every cardinality, and the location of the de re/de dicto
separation --- are ones we have not found stated, but they are elementary
consequences of results in the cited work, and we do not claim them as new.
Wernhard~\cite{wernhard} computes a frame condition for Scott's Theorem 3 by
second-order quantifier elimination and leaves open whether it is the weakest;
that question is the nearest thing to an opening this material still has.

\section{The system}\label{sec:system}

\subsection{Semantics}

We work with a shallow embedding of quantified modal logic, which is the
standard device for this kind of mechanization and is what makes the
formalization short.

Fix a type $W$ of possible worlds and a type $I$ of individuals. A
\emph{sentence} is a map $W \to \Prp$, and a \emph{property} is a map
$I \to W \to \Prp$; a property thus has an extension that may vary from world
to world. Let $r \subseteq W \times W$ be an accessibility relation. We write
\[
  (\Box_r p)(w) \;:=\; \forall v\,\bigl(r\,w\,v \to p(v)\bigr),
  \qquad
  (\Diamond_r p)(w) \;:=\; \exists v\,\bigl(r\,w\,v \wedge p(v)\bigr).
\]
Quantifiers over individuals and over properties are read world by world, with
a constant domain. Quantification over properties ranges over \emph{all} maps
$I \to W \to \Prp$; see Section~\ref{sec:caveats} for why this matters.

\emph{No condition is imposed on $r$ as part of the semantics.} Where a
derivation needs one it is stated as a hypothesis of that derivation and
nowhere else. This is the whole reason for keeping $r$ explicit rather than
fixing it to the universal relation, and Section~\ref{sec:frames} collects what
it buys.

An axiom is asserted to hold \emph{at every world}. This is the ordinary
reading, and Section~\ref{sec:incons} uses it in an essential way.

\subsection{Definitions}

The sole nonlogical primitive is
$\Pos : (I \to W \to \Prp) \to (W \to \Prp)$, read ``$\varphi$ is a positive
property'', and itself world-dependent: Axiom~4 asserts
$\Pos(\varphi) \to \Box \Pos(\varphi)$, which would be idle if $\Pos$ were not
allowed to vary.

\begin{definition}[God-likeness]
$\Gd(x) \;:=\; \forall \varphi\,\bigl(\Pos(\varphi) \to \varphi(x)\bigr)$.
\end{definition}

\begin{definition}[Essence]\label{def:ess}
G\"odel's 1970 definition and Scott's 1987 emendation are
\begin{align*}
  \varphi \ess x  &\;:=\;
    \forall \psi\,\Bigl(\psi(x) \to \Box\,\forall y\,\bigl(\varphi(y) \to \psi(y)\bigr)\Bigr), \\[2pt]
  \varphi \esss x &\;:=\;
    \varphi(x) \;\wedge\;
    \forall \psi\,\Bigl(\psi(x) \to \Box\,\forall y\,\bigl(\varphi(y) \to \psi(y)\bigr)\Bigr).
\end{align*}
\end{definition}

\noindent
Both read: everything $x$ has, $\varphi$ necessarily carries with it. Scott
adds the requirement that $x$ actually \emph{have} $\varphi$. That conjunct is
the entire difference between the two systems, and everything below turns on
it.

\begin{definition}[Necessary existence]
$\NE(x) := \forall \varphi\,(\varphi \ess x \to \Box\,\exists y\,\varphi(y))$,
and $\NES(x)$ is the same with $\esss$ in place of $\ess$.
\end{definition}

\subsection{The axioms}

\begin{description}
  \item[A1a.] $\Pos(\neg\varphi) \to \neg\,\Pos(\varphi)$.
  \item[A1b.] $\neg\,\Pos(\varphi) \to \Pos(\neg\varphi)$.
  \item[A2.] $\Pos(\varphi) \wedge \Box\,\forall x\,(\varphi(x) \to \psi(x))
             \;\to\; \Pos(\psi)$.
  \item[A3.] $\Pos(\Gd)$.
  \item[A4.] $\Pos(\varphi) \to \Box\,\Pos(\varphi)$.
  \item[A5.] $\Pos(\NE)$ \quad(G\"odel) \qquad or \qquad $\Pos(\NES)$ \quad(Scott).
\end{description}

\noindent
A1 is stated by both authors as the biconditional; we split it because the two
halves are used by different results. \emph{G\"odel's system} is
A1a, A1b, A2, A3, A4 together with $\Pos(\NE)$; \emph{Scott's system} is the
same with $\Pos(\NES)$.

\section{G\"odel's 1970 axioms are inconsistent}\label{sec:incons}

The derivation has two moving parts. The first is G\"odel's own Theorem~1, and
is uncontroversial; it is shared by both systems.

\begin{lemma}[Theorem 1]\label{lem:T1}
Assume {\rm A1a} and {\rm A2}. If $\Pos(\varphi)$ holds at $w$, then
$(\Diamond_r \exists x\, \varphi(x))(w)$.
\end{lemma}

\begin{proof}
Suppose not, so that $\varphi$ has empty extension at every world accessible
from $w$. Let $\psi$ be arbitrary. Then
$\Box_r \forall x\,(\varphi(x) \to \psi(x))$ holds at $w$ vacuously: at any
accessible $v$ and any $y$, the antecedent $\varphi(y)$ fails. By A2 together
with $\Pos(\varphi)$ we get $\Pos(\psi)$ at $w$. As $\psi$ was arbitrary,
\emph{every} property is positive at $w$; in particular both $\Pos(\neg\varphi)$
and $\Pos(\varphi)$, contradicting A1a.
\end{proof}

The second part is where the 1970 system comes apart.

\begin{lemma}[Empty Essence Lemma~\cite{bwp2016}]\label{lem:empty}
Let $e$ be the property with empty extension at every world. Then for every
individual $x$ and every world $w$, $\,e \ess x\,$ holds at $w$.
\end{lemma}

\begin{proof}
Let $\psi$ be any property with $\psi(x)$ at $w$; we must show
$\Box_r \forall y\,(e(y) \to \psi(y))$ at $w$. Let $v$ be accessible from $w$
and $y$ any individual. The antecedent $e(y)$ is false at $v$ by construction,
so the implication holds.
\end{proof}

\noindent
No axiom, no frame condition and no classical reasoning is used here: the
statement is a triviality about vacuous implication. That is precisely what
makes it dangerous. Definition~\ref{def:ess} places a \emph{demand} on
$\varphi$ --- that it necessarily entail everything $x$ has --- and a property
that nothing satisfies meets every such demand for free.

\begin{remark}
Benzm\"uller and Woltzenlogel Paleo observe that some philosophers may prefer
to run the argument with the property of \emph{self-difference},
$\lambda x.\, x \neq x$, in place of the empty property. Under full
extensionality the two have the same denotation and the derivation is
indifferent to which is used.
\end{remark}

\begin{theorem}[\cite{bwp2016}]\label{thm:main}
Suppose $W$ is nonempty. Then {\rm A1a}, {\rm A2} and $\Pos(\NE)$ are jointly
unsatisfiable, over an arbitrary accessibility relation $r$.
\end{theorem}

\begin{proof}
By A5, $\Pos(\NE)$ holds at every world, so Lemma~\ref{lem:T1} gives
\begin{equation}\label{eq:star}
  \text{for every } w \in W \text{ there are } v \text{ with } r\,w\,v
  \text{ and } x \in I \text{ with } \NE(x) \text{ at } v. \tag{$*$}
\end{equation}

Pick any $w_0 \in W$ and apply \eqref{eq:star}, obtaining such a $v$ and $x$.
Unfolding, $\NE(x)$ at $v$ says
$\forall \varphi\,(\varphi \ess x \to \Box_r \exists y\, \varphi(y))$.
Instantiate at $\varphi := e$. By Lemma~\ref{lem:empty} the antecedent holds,
so $\Box_r \exists y\, e(y)$ holds at $v$: at every world accessible from $v$,
something has the empty property. Nothing has the empty property at any world.
Hence
\begin{equation}\label{eq:dead}
  \text{no world is accessible from } v. \tag{$\dagger$}
\end{equation}

But \eqref{eq:star} holds at \emph{every} world, $v$ included, because the
axioms do. Applying it at $v$ produces a world accessible from $v$, in direct
contradiction with~\eqref{eq:dead}.
\end{proof}

\noindent
Note that neither A3 nor A4 is used: the argument never needs a God-like being
to exist, because Lemma~\ref{lem:T1} is applied to $\NE$ directly. The last
step is what keeps the frame conditions at zero. A cruder argument would cash
$\Box_r \exists y\, e(y)$ into a witness by assuming $r$ serial, at the cost of
working in $D$; that is unnecessary, since the axioms are asserted everywhere
and Lemma~\ref{lem:T1} can simply be reapplied at $v$.

\begin{proposition}\label{prop:scott}
For every individual $x$ and every world $w$, $\,e \esss x\,$ fails at $w$.
\end{proposition}

\begin{proof}
It would require $e(x)$ at $w$, and $e$ has empty extension.
\end{proof}

\noindent
So the Empty Essence Lemma --- the only input to Theorem~\ref{thm:main} beyond
Theorem~1 --- is unavailable in Scott's system, and the derivation stops. This
does not by itself show Scott's system consistent; that is
Proposition~\ref{prop:scottcons} below.

\section{Scott's variant: the positive argument}\label{sec:scott}

We now carry out the argument G\"odel intended. Throughout this section the
axioms are Scott's.

\begin{lemma}\label{lem:posgod}
Assume {\rm A1b}. If $\Gd(x)$ and $\psi(x)$ both hold at $w$, then
$\Pos(\psi)$ holds at $w$.
\end{lemma}

\begin{proof}
If $\Pos(\psi)$ failed at $w$ then $\Pos(\neg\psi)$ would hold there by A1b;
since $x$ is God-like it has every positive property, so $\neg\psi(x)$ at $w$,
contradicting $\psi(x)$.
\end{proof}

\begin{theorem}[Theorem 2]\label{thm:T2}
Assume {\rm A1b} and {\rm A4}. If $\Gd(x)$ holds at $w$, then
$\Gd \esss x$ holds at $w$. No frame condition is needed.
\end{theorem}

\begin{proof}
The first conjunct of $\esss$ is $\Gd(x)$ at $w$, which is the hypothesis. For
the second, let $\psi$ satisfy $\psi(x)$ at $w$. By Lemma~\ref{lem:posgod},
$\Pos(\psi)$ at $w$, so by A4, $\Pos(\psi)$ holds at every world $v$ accessible
from $w$. Let $v$ be such a world and let $y$ satisfy $\Gd(y)$ at $v$. Being
God-like, $y$ has every property positive at $v$, and $\psi$ is one; so
$\psi(y)$ at $v$. This is exactly
$\Box_r \forall y\,(\Gd(y) \to \psi(y))$ at $w$.
\end{proof}

\begin{lemma}\label{lem:boxgod}
Assume {\rm A1b}, {\rm A4} and $\Pos(\NES)$. If $\Gd(g)$ holds at $w$, then
$\Box_r \exists y\, \Gd(y)$ holds at $w$.
\end{lemma}

\begin{proof}
$\Pos(\NES)$ at $w$ and $\Gd(g)$ at $w$ give $\NES(g)$ at $w$. By
Theorem~\ref{thm:T2}, $\Gd \esss g$ at $w$. Instantiate $\NES(g)$ at
$\varphi := \Gd$.
\end{proof}

\noindent
This is the step Scott's conjunct rescues. In G\"odel's system the same move,
made with $\varphi := e$, is what produced Theorem~\ref{thm:main}.

\begin{lemma}\label{lem:existsgod}
Assume Scott's axioms and that $r$ is \emph{symmetric}. Then a God-like
individual exists at every world.
\end{lemma}

\begin{proof}
Let $w$ be any world. By A3 and Lemma~\ref{lem:T1} there are $v$ with $r\,w\,v$
and $g$ with $\Gd(g)$ at $v$. By Lemma~\ref{lem:boxgod},
$\Box_r \exists y\, \Gd(y)$ holds at $v$. By symmetry $r\,v\,w$, so
$\exists y\, \Gd(y)$ holds at $w$.
\end{proof}

\begin{theorem}[Theorem 3]\label{thm:T3}
Assume Scott's axioms and that $r$ is symmetric. Then
$\Box_r \exists x\, \Gd(x)$ holds at every world.
\end{theorem}

\begin{proof}
Combine Lemma~\ref{lem:existsgod} with Lemma~\ref{lem:boxgod}.
\end{proof}

\noindent
This is the conclusion of the ontological argument, and it is worth being
plain about what it is. It is a \emph{conditional}: it says that granting
Scott's five axioms commits one to $\Box\exists x\,\Gd(x)$. Whether they should
be granted is not a question a proof assistant can address. What a proof
assistant \emph{can} do is compute the rest of the bill, which is the subject
of the next section.

\subsection{A shorter route, and a warning about it}\label{sec:shortcut}

The derivation above is Scott's, and it uses A1b and A4 through
Theorem~\ref{thm:T2}. Our encoding also admits a route that uses neither. We
record it because it is a genuine theorem of the encoding and because the gap
between it and Scott's argument is itself informative --- but the warning
attached to it is not a formality, and we state it before the lemma rather
than after.

\begin{lemma}[Haecceities are essences]\label{lem:haecceity}
Fix $g$ and $v$ and let $\varphi := \lambda z\,u.\,(z = g \wedge u = v)$. Then
$\varphi \esss g$ holds at $v$, with no axioms and over any frame.
\end{lemma}

\begin{proof}
The first conjunct is $g = g \wedge v = v$. For the second, let $\psi$ be a
property with $\psi(g)$ at $v$, let $r\,v\,u$, and let $y$ satisfy
$\varphi(y)$ at $u$. Then $y = g$ and $u = v$, so what is required is
$\psi(g)$ at $v$, which is the hypothesis.
\end{proof}

\begin{lemma}\label{lem:blind}
If $\Gd(g)$ at $v$, then A5 alone gives $r\,v\,u \to u = v$ for all $u$.
\end{lemma}

\begin{proof}
By D1 and A5, $\NES(g)$ holds at $v$. Apply it to the $\varphi$ of
Lemma~\ref{lem:haecceity}: it yields $\Box_r \exists y\, \varphi(y)$ at $v$,
and the only way $\varphi(y)$ can hold at $u$ is for $u$ to be $v$.
\end{proof}

\noindent
Lemma~\ref{lem:blind} does on its own what Theorem~\ref{thm:frame} does with
the full axiom set: A5 together with the definition of God-likeness seals off
the world of a God-like individual. Theorem~\ref{thm:T3}, modal collapse and
Theorem~\ref{thm:frame} all follow from it together with A1a, A2, A3 and
symmetry, and our Lean statements of those results now carry exactly those
hypotheses.

\paragraph{The status of Lemma~\ref{lem:haecceity}.} The reduction depends on
the property domain containing world-indexed haecceities such as
$H_{g,v}(z,u) :\equiv (z = g \wedge u = v)$. Our shallow embedding contains
every function $I \to W \to \Prp$ and so contains this one. Whether that is
faithful to Scott's intended quantification over properties depends on the
semantics adopted, and there are two regimes.

Under \emph{full intensional semantics}, where the property domain really is
all of $I \times W \to \Prp$, $H_{g,v}$ is a perfectly legitimate property and
the shortcut is a genuine consequence of Scott's definitions together with A5.
Under a \emph{restricted or Henkin} property domain --- or one given by a modal
comprehension schema over the object language --- the question is whether
$H_{g,v}$ belongs to the domain at all. A standard modal object language has
no term naming the current world, so definability is not automatic, and the
shortcut may fail. Some presentations of Scott state the background theory
with a comprehension schema rather than identifying properties with arbitrary
sets of individual--world pairs, which is precisely why the distinction bites.

One concrete data point. The Isabelle/HOL embedding used by Benzm\"uller and
collaborators (AFP, \emph{Types, Tableaus and Gödel's God}, theory
\verb|IHOML|) types a property as $\mathsf{ise} = \mathbf{0} \Rightarrow
(i \Rightarrow \mathrm{bool})$ --- the same function space as ours --- and
quantifies over it with plain HOL quantification, with no comprehension schema
restricting property formation. So $H_{g,v}$ is expressible and in range there
as well, and the reduction is unlikely to be peculiar to Lean. That is an
inspection of their type declarations, not a proof run: we have not asked their
development for Lemma~\ref{lem:blind}, and doing so is the natural next
experiment. We make \emph{no claim of novelty} either way.

The lemma is reported because a formalization that silently admitted a
shortcut it did not mention would be the worse for it, and because locating
the exact semantic assumption a reduction rests on is one of the things
formalization is for.

Fitting's variant makes the contrast sharp. The same idea works there with no
warning at all: his essences are extensions, so the bare singleton $\{g\}$
suffices and no world is named (Lemma~\ref{lem:singleton}). That the move is
clean under the extensional reading and suspect under the intensional one is
the intension/extension distinction reappearing as a difference in what the
essence clause permits.

\section{The price: modal collapse}\label{sec:collapse}

\begin{theorem}[Modal collapse, Sobel]\label{thm:collapse}
Assume Scott's axioms and that $r$ is symmetric. Then for every sentence $p$
and every world $w$, if $p$ holds at $w$ then $\Box_r p$ holds at $w$.
\end{theorem}

\begin{proof}
By Lemma~\ref{lem:existsgod} there is $g$ with $\Gd(g)$ at $w$, and by
Theorem~\ref{thm:T2}, $\Gd \esss g$ at $w$. Regard $p$ as a property,
$\psi_p := \lambda y.\, p$, constant in the individual. Then $\psi_p(g)$ at $w$
is just $p$ at $w$, which holds; so the second conjunct of $\esss$ gives
\[
  \Box_r \forall y\,\bigl(\Gd(y) \to p\bigr) \quad\text{at } w .
\]
By Theorem~\ref{thm:T3}, $\Box_r \exists x\,\Gd(x)$ at $w$. So at each world
accessible from $w$ there is a God-like individual, and hence $p$ holds there.
\end{proof}

\noindent
The proof just given is Scott's. Our Lean \verb|modal_collapse| takes the
shorter route of \S\ref{sec:shortcut} instead and so consumes only A1a, A2, A3,
A5 and symmetry. The two are worth keeping apart, because "what does collapse
need?" has different answers depending on whether Theorem 3 is expanded or
taken as given. The relative question has a clean answer, and it is the one the
published summaries record.

\begin{proposition}[Collapse relative to Theorem 3]\label{thm:mcrel}
Assume $r$ symmetric and serial, assume A1b and A4, and assume
$\Box_r \exists x\,\Gd(x)$ at every world. Then modal collapse holds. None of
A1a, A2, A3, A5 is used.
\end{proposition}

\begin{proof}
Seriality gives $v$ with $r\,w\,v$; symmetry gives $r\,v\,w$; the hypothesis at
$v$ then puts a God-like $g$ at $w$. Theorem~\ref{thm:T2}, which needs only A1b
and A4, makes $\Gd \esss g$ at $w$, and the argument proceeds as above.
\end{proof}

\begin{corollary}
If $r$ is in addition reflexive, then $p \leftrightarrow \Box_r p$ for every
sentence $p$: the distinction between truth and necessary truth disappears.
\end{corollary}

Nothing is contingent. This is Sobel's objection, and it is worth stressing
that it applies to Scott's \emph{consistent} system: it is not a symptom of the
1970 defect, and repairing that defect does not touch it. Under the $S5$
reading the consequence can be put more starkly still.

\begin{theorem}\label{thm:oneworld}
Assume Scott's axioms with $r$ the universal relation. Then any two worlds are
equal: $W$ has at most one element.
\end{theorem}

\begin{proof}
Apply Theorem~\ref{thm:collapse} to the sentence $p := \lambda u.\, u = w$,
which holds at $w$. Then $\Box_r p$ at $w$, and since every world is accessible,
$v = w$ for every $v$.
\end{proof}

\subsection{Collapse as a condition on the frame}\label{sec:frame}

The same instance of Theorem~\ref{thm:collapse} says something more general,
and more informative, once $r$ is not assumed universal. Modal collapse is
usually stated as a fact about sentences; its semantic content is a fact about
the accessibility relation.

\begin{theorem}\label{thm:frame}
Assume Scott's axioms and that $r$ is symmetric. Then $r$ is \emph{equality}:
$r\,w\,v$ holds if and only if $v = w$.
\end{theorem}

\begin{proof}
For one direction, apply Theorem~\ref{thm:collapse} at $w$ to
$p := \lambda u.\, u = w$, which holds at $w$; this gives $\Box_r p$ at $w$,
that is, $r\,w\,v \to v = w$. For the other, Lemma~\ref{lem:T1} applied to A3
produces some $v$ with $r\,w\,v$; by what we have just shown $v = w$, so
$r\,w\,w$.
\end{proof}

\noindent
Two features of this deserve comment. The frame conditions are \emph{derived},
not assumed: reflexivity and seriality are consequences of the axioms, not
hypotheses added to make the argument work. And the conclusion is not that
there are few worlds, but that no world can see another.

\begin{corollary}\label{cor:inert}
Under the same hypotheses, $\Box_r p \leftrightarrow p$ and
$\Diamond_r p \leftrightarrow p$ for every sentence $p$.
\end{corollary}

So the modal operators are the identity operation. Theorem~\ref{thm:oneworld}
is the special case in which $r$ is additionally required to be universal:
equality and the universal relation agree only on a one-point set. But the
general statement is the one that bites. The argument is conducted in a modal
language about what is possible and what is necessary, and its own premises
leave that language with nothing to say: every world is an isolated point,
$\Box$ means nothing beyond ``here'', and the conclusion
$\Box\exists x\,\Gd(x)$ reduces to $\exists x\,\Gd(x)$.

\begin{remark}
This is the standard semantic reading of modal collapse rather than a new
observation; the condition $\forall x\,\forall y\,(r\,x\,y \to y = x)$ is
stated directly as a rendering of $MC$ in the literature on computational
variants of the argument. What Theorem~\ref{thm:frame} adds here is that it is
\emph{derived} from Scott's axioms in the formalization, with the frame
hypothesis (symmetry, and nothing else) recorded.
\end{remark}

\section{One model, two jobs}\label{sec:model}

Let $W = I = \{*\}$, let $r$ be the universal relation, and put
$\Pos(\varphi) :\Leftrightarrow \varphi(*)$ --- a property is positive exactly
when the sole individual has it.

\begin{proposition}\label{prop:toy}
In this model {\rm A1a}, {\rm A1b}, {\rm A2}, {\rm A3} and {\rm A4} all hold,
and G\"odel's {\rm A5}, $\Pos(\NE)$, fails.
\end{proposition}

\begin{proof}
A1a and A1b both read $\neg\varphi(*) \to \neg\varphi(*)$. For A2, suppose
$\varphi(*)$ and $\Box\forall x(\varphi(x)\to\psi(x))$; instantiating at the
unique world and individual gives $\psi(*)$. For A3 we must check $\Gd(*)$,
that is $\forall\varphi(\Pos(\varphi) \to \varphi(*))$, which unfolds to
$\forall\varphi(\varphi(*) \to \varphi(*))$. A4 is immediate, there being one
world. Finally $\Pos(\NE)$ fails: it would give $\NE(*)$, and instantiating at
$e$, which is a G\"odel-essence of $*$ by Lemma~\ref{lem:empty}, yields
$\exists y\, e(y)$.
\end{proof}

This is the non-triviality check that Theorem~\ref{thm:main} needs. A
mechanized proof of $\bot$ from a list of hypotheses is only informative if the
hypotheses have not been mis-stated into absurdity; nothing in
Theorem~\ref{thm:main} would look different if A2 had been encoded in a form
nothing could satisfy. Proposition~\ref{prop:toy} rules that out and isolates
A5 as the axiom that tips the system over.

The same model then settles a second question.

\begin{proposition}\label{prop:scottcons}
In the same model, with the same $\Pos$, Scott's {\rm A5}, $\Pos(\NES)$, holds.
Hence Scott's system is consistent.
\end{proposition}

\begin{proof}
We must check $\NES(*)$, that is, that every Scott-essence of $*$ is
necessarily instantiated. Let $\varphi \esss *$. The first conjunct gives
$\varphi(*)$, and with one world and one individual that is already
$\Box\exists y\, \varphi(y)$.
\end{proof}

The contrast between Propositions~\ref{prop:toy} and~\ref{prop:scottcons} is
the cleanest available statement of what G\"odel's omission costs. One
structure, one notion of positivity; G\"odel's Axiom~5 fails in it and Scott's
holds, and the only difference between them is the conjunct $\varphi(x)$. It
also shows that the inconsistency of Section~\ref{sec:incons} is specific to
the definition of essence rather than a defect of the surrounding axioms or of
our encoding of them.

This model is degenerate: one world, so $\Box$ is trivial and modal collapse
holds vacuously. It is natural to ask whether that is an accident of choosing
the smallest structure --- whether Scott's axioms have a model with several
worlds standing in some genuine accessibility relation. Theorem~\ref{thm:frame}
already answers the second half, but the first half deserves an explicit
answer, since a theorem forbidding modal structure would be idle if it also
forbade models outright.

\begin{proposition}\label{prop:iso}
Let $V$ be any type of worlds whatsoever, let $r$ be equality on $V$, take one
individual, and put $\Pos(\varphi)$ at $w$ iff that individual has $\varphi$ at
$w$. Then all of Scott's axioms hold.
\end{proposition}

\begin{proof}
Each world sees only itself, so $\Box_r q$ at $w$ is just $q$ at $w$, and every
axiom reduces to its one-world form, verified as in
Proposition~\ref{prop:scottcons}.
\end{proof}

\noindent
So the answer to ``does Scott's system have a non-degenerate model?'' is:
\emph{yes in cardinality, no in modal structure}. The world set may be as large
as one likes --- Proposition~\ref{prop:iso} places no hypothesis on $V$ at all
--- but by Theorem~\ref{thm:frame} every such model is a disjoint heap of
isolated copies of Proposition~\ref{prop:scottcons}, each with its own God-like
individual, none able to see any other. The one-world model is not forced; but
nothing better than a heap of one-world models is available.

\section{Symmetry cannot be traded for $S4$}\label{sec:symm}

Symmetry is used in exactly one place, Lemma~\ref{lem:existsgod}, to carry
$\Box\exists x\,\Gd(x)$ from a merely possible world back to the actual one.
That is the $B$ schema $\Diamond\Box p \to p$. Everything downstream ---
Theorem~\ref{thm:T3}, modal collapse, Theorem~\ref{thm:frame} --- inherits the
hypothesis from there. It is natural to ask whether it can be weakened.

It cannot be traded for the other two $S5$ conditions: reflexivity and
transitivity are no substitute. Note what this does and does not show. It
establishes that the symmetry hypothesis cannot simply be deleted from the
$KB$ result in favour of the $S4$ conditions. It does \emph{not} establish
that symmetry is necessary in the absolute sense; some frame condition weaker
than, or incomparable with, symmetry might still suffice, and we prove no
correspondence or minimality theorem that would rule this out.

\begin{proposition}\label{prop:nosymm}
There is a model of Scott's axioms over a reflexive, transitive, non-symmetric
frame in which Theorem~\ref{thm:T3}, Theorem~\ref{thm:collapse} and
Theorem~\ref{thm:frame} all fail.
\end{proposition}

\begin{proof}
Let $W = \{0,1\}$ with $r\,w\,v$ iff $w = 0$ or $v = 1$: world $0$ sees both
worlds, world $1$ sees only itself. This is reflexive and transitive, and not
symmetric, since $r\,0\,1$ but not $r\,1\,0$. Take one individual $*$, and put
\[
  \Pos(\varphi)(w) \;:\Leftrightarrow\; \varphi(*)(1)
  \qquad\text{at both worlds}.
\]
So positivity is imported wholesale from world $1$, where the model of
Proposition~\ref{prop:scottcons} lives.

A1a and A1b are immediate. For A2, the box hypothesis at either world
instantiates at $v = 1$, which is accessible from both, and the argument of
Proposition~\ref{prop:scottcons} applies. A3 asks for $\Gd(*)$ at $1$, which
unfolds to $\forall\varphi(\varphi(*)(1) \to \varphi(*)(1))$. A4 holds because
$\Pos(\varphi)$ is a constant sentence. A5 asks for $\NES(*)$ at $1$; every
world accessible from $1$ is $1$ itself, so the demand reduces to the first
conjunct of $\esss$, exactly as before.

Now $\Gd(*)$ \emph{fails} at world $0$. Being God-like at $0$ would require
having at $0$ every property positive there, hence every property one has at
$1$; and $\lambda y\,w.\,(w = 1)$ is such a property, false at $0$. So there is
no God-like individual at $0$. Since $r\,0\,0$, that refutes
$\Box\exists x\,\Gd(x)$ at $0$, which is Theorem~\ref{thm:T3}. Modal collapse
fails at $0$ as well: $\lambda u.\,(u=0)$ holds at $0$ but not at the
accessible world $1$. And $r\,0\,1$ with $1 \neq 0$ refutes
Theorem~\ref{thm:frame}.
\end{proof}

Two things follow. First, symmetry is not a convenience of the usual $S5$
presentation that a sharper proof could remove; the $B$ schema is doing
irreplaceable work, and it is doing it at precisely the step where the argument
moves from possible existence to actual existence. Second, and more striking:

\begin{corollary}\label{cor:nocollapse}
Scott's axioms do not entail modal collapse.
\end{corollary}

\noindent
They entail it \emph{in $B$ and $S5$} --- which is where the argument is always
presented, so the objection stands where it is usually made. But collapse is a
consequence of the axioms \emph{together with} a symmetric frame, not of the
axioms alone, and Proposition~\ref{prop:nosymm} exhibits an $S4$ model of the
axioms in which contingency survives. The price of the argument is therefore
better stated as a price of the frame it is run in: strengthen the logic enough
to deliver the conclusion, and you get the collapse with it.

\begin{remark}
This is consistent with, and gives one explanation of, the report
in~\cite{bwp2016} that the argument needs $KB$ --- symmetry and nothing more.
It also locates the target for the emendations, which is the subject of
Section~\ref{sec:variants}.
\end{remark}

\section{The emendations}\label{sec:variants}

Anderson~\cite{anderson} and Fitting~\cite{fitting} each modify the argument so
as to keep the conclusion and lose the collapse. Benzm\"uller and
Fuenmayor~\cite{bf2020} have analysed both on the computer; we follow their
presentation, and in particular their rendering of Anderson's variant as
Fitting gives it. This section formalizes each variant far enough to prove its
Theorem 3 and to exhibit a model in which modal collapse fails.

\subsection{Anderson: biconditionals in place of {\rm A1b}}

Anderson's complaint is about A1b, $\neg\Pos(\varphi) \to \Pos(\neg\varphi)$,
which says that of every property one of it and its negation is positive. That
is a strong claim, and --- as Section~\ref{sec:incons} showed --- not one the
refutation of G\"odel's system ever needed. Anderson drops it. Dropping it
alone breaks Theorem 2, so he strengthens the two definitions to compensate:
\begin{align*}
  \Gd^{\mathrm{A}}(x) &\;:=\; \forall\varphi\,\bigl(\Pos(\varphi) \leftrightarrow \Box\varphi(x)\bigr), \\
  \varphi \ess^{\mathrm{A}} x &\;:=\; \forall\psi\,\bigl(\Box\psi(x) \leftrightarrow \Box\forall y(\varphi(y)\to\psi(y))\bigr),
\end{align*}
with $\NE^{\mathrm{A}}$ as before but using $\ess^{\mathrm{A}}$. The axioms are
A1a, A2, A4, $\Pos(\NE^{\mathrm{A}})$, and $\Pos(\Gd^{\mathrm{A}})$ postulated
in place of A3.

The biconditional in $\Gd^{\mathrm{A}}$ does far more work than it first
appears to, and the following one-line lemma is the reason. It is the engine of
the entire variant.

\begin{lemma}[God-likeness is self-propagating]\label{lem:andboxgod}
From $\Pos(\Gd^{\mathrm{A}})$ alone: if $\Gd^{\mathrm{A}}(x)$ at $w$, then
$\Box\Gd^{\mathrm{A}}(x)$ at $w$. No frame condition, and no other axiom.
\end{lemma}

\begin{proof}
Instantiate the biconditional $\Pos(\varphi) \leftrightarrow \Box\varphi(x)$ of
$\Gd^{\mathrm{A}}(x)$ at $\varphi := \Gd^{\mathrm{A}}$. The left-hand side is
Anderson's A3$^{*}$.
\end{proof}

\noindent
Where G\"odel and Scott reach the corresponding step through the whole essence
apparatus together with A5, Anderson gets it by feeding a definition to itself.
Uniqueness comes almost as cheaply.

\begin{proposition}[Monotheism]\label{prop:mono}
Anderson's axioms make God-like individuals unique, over \emph{any} frame: if
$\Gd^{\mathrm{A}}(x)$ and $\Gd^{\mathrm{A}}(y)$ at $w$, then $y = x$.
\end{proposition}

\begin{proof}
Let $\varphi := \lambda z\,u.\,(z = x)$. Then $\Box\varphi(x)$ holds at $w$, so
$\Pos(\varphi)$ at $w$ by $\Gd^{\mathrm{A}}(x)$; so $\Box\varphi(y)$ at $w$ by
$\Gd^{\mathrm{A}}(y)$. Theorem~\ref{lem:T1} applied to $\Gd^{\mathrm{A}}$,
which is positive by A3$^{*}$, supplies a world $v$ with $r\,w\,v$; evaluating
$\Box\varphi(y)$ there gives $\varphi(y)$, i.e.\ $y = x$.
\end{proof}

\noindent
An earlier version of this note assumed reflexivity for that last step. It is
not needed: A1a, A2 and A3$^{*}$ already guarantee that $w$ has a successor,
which is all the argument wants.

\noindent
Uniqueness is \emph{not} peculiar to Anderson, and an earlier version of this
note said it was. The same argument runs in Scott's system, from A1b alone and
with no frame condition at all: if being identical to a God-like $x$ were not
positive then its negation would be, by A1b, and $x$ --- having every positive
property --- would fail to be self-identical. So it is positive, and any
God-like $y$ has it. This is \verb|Godel.scott_monotheism| in the development,
and Benzm\"uller and Scott~\cite{bs2025} likewise record monotheism as a
consequence of the G\"odel/Scott setting. What Anderson's biconditional buys is
directness, not the theorem.

\begin{theorem}\label{thm:andT3}
Over a \emph{symmetric} frame --- the logic $KB$ --- A1a, A2 and
$\Pos(\Gd^{\mathrm{A}})$ entail $\Box\exists x\,\Gd^{\mathrm{A}}(x)$. Neither
A4 nor A5 is used.
\end{theorem}

\begin{proof}
$\Gd^{\mathrm{A}}$ is positive by A3$^{*}$, so Theorem~\ref{lem:T1} --- which
needs only A1a and A2 --- gives a world $v$ with $r\,w\,v$ and an $x$ with
$\Gd^{\mathrm{A}}(x)$ at $v$. By Lemma~\ref{lem:andboxgod},
$\Box\Gd^{\mathrm{A}}(x)$ at $v$; symmetry gives $r\,v\,w$, so
$\Gd^{\mathrm{A}}(x)$ at $w$. Applying Lemma~\ref{lem:andboxgod} again at $w$
gives $\Box\Gd^{\mathrm{A}}(x)$ at $w$, and in particular
$\Box\exists x\,\Gd^{\mathrm{A}}(x)$.
\end{proof}

\noindent
A1a and A2 enter only through Theorem~\ref{lem:T1}, and only to guarantee that
$w$ has a successor at all. Symmetry is used once, to bring the individual home
from the world where Theorem~1 found it.

This matches what is already known. Anderson himself (1990, fn.~5) notes that
the weaker logic $B$ --- that is, $KTB$ --- suffices for Theorem 3, and
Benzm\"uller, Weber and Woltzenlogel Paleo~\cite{bww2017} report it automated
in the weaker $KB$ again. \emph{We claim no new mathematics here.} What is new
in this development is only that the frame condition and the axiom list are now
carried by the type of Theorem~\ref{thm:andT3} --- it takes a record containing
A1a, A2 and A3$^{*}$ and nothing else --- rather than being recorded in prose
alongside a proof that in fact used more. An earlier version of this note
derived Theorem~\ref{thm:andT3} over $S5$ from the full axiom set and flagged
the gap to $KB$ as open work; that route survives as
\verb|Anderson.T3_viaEssence|, and the contrast is instructive, since it shows
the $S5$ was a property of the argument rather than of Anderson's system.

\subsubsection*{A4 and A5 are redundant}

The same paper reports A4 and A5 \emph{redundant} for this emendation, A5
already in $K$ and A4 in $K4B$. Both are now proved rather than cited. The
first is stronger than redundancy.

\begin{proposition}\label{prop:nevac}
With Anderson's essence, $\NE^{\mathrm{A}}(x)$ holds for every $x$ at every
world, with no axioms and over any frame.
\end{proposition}

\begin{proof}
Let $\varphi \ess^{\mathrm{A}} x$ at $w$ and instantiate its biconditional at
$\psi := \varphi$. The right-hand side is
$\Box\forall y\,(\varphi(y) \to \varphi(y))$, a validity, so $\Box\varphi(x)$
holds at $w$. Hence at every $v$ with $r\,w\,v$ we have $\varphi(x)$, so
$\exists y\,\varphi(y)$; that is, $\Box\exists y\,\varphi(y)$ at $w$.
\end{proof}

\begin{corollary}\label{cor:a5}
A2 and A3$^{*}$ entail $\Pos(\NE^{\mathrm{A}})$, over any frame.
\end{corollary}

\begin{proof}
$\Pos(\Gd^{\mathrm{A}})$ by A3$^{*}$, and
$\Box\forall x\,(\Gd^{\mathrm{A}}(x) \to \NE^{\mathrm{A}}(x))$ holds because
its consequent is true outright, by Proposition~\ref{prop:nevac}. Apply A2.
\end{proof}

\noindent
So A5 is not merely derivable in Anderson's system: it is contentless there,
because $\NE^{\mathrm{A}}$ is the universally true property. The axiom
philosophers argue over hardest in G\"odel's system does no work at all in
Anderson's.

\begin{proposition}\label{prop:a4}
Over a symmetric and transitive frame --- the logic $K4B$ --- A1a, A2 and
A3$^{*}$ entail A4.
\end{proposition}

\begin{proof}
Let $\Pos(\varphi)$ at $w$ and $r\,w\,u$. As in Theorem~\ref{thm:andT3} there
is a $g$ with $\Gd^{\mathrm{A}}(g)$ at $w$, and by Lemma~\ref{lem:andboxgod} also
at $u$. From $\Gd^{\mathrm{A}}(g)$ at $w$ and $\Pos(\varphi)$ at $w$ we get
$\Box\varphi(g)$ at $w$. For any $u'$ with $r\,u\,u'$, transitivity gives
$r\,w\,u'$ and hence $\varphi(g)$ at $u'$; so $\Box\varphi(g)$ at $u$, and
$\Gd^{\mathrm{A}}(g)$ at $u$ converts that to $\Pos(\varphi)$ at $u$.
\end{proof}

\noindent
Both frame conditions agree with those reported in~\cite{bww2017}.

\paragraph{Scope.} All of the above is for the constant-domain (possibilist)
reading used throughout this note. Under the \emph{mixed} reading Anderson
himself floats --- actualist quantifiers only in Theorem 3 and in the
definition of essence ---~\cite{bww2017} find A4 still redundant but A5
\emph{independent}, and report a countermodel to Theorem 3 itself. Domain
choice is not a free parameter, and none of Theorem~\ref{thm:andT3},
Corollary~\ref{cor:a5} or Proposition~\ref{prop:a4} transfers there.

\subsubsection*{Anderson's own route}

The essence argument is retained in the development for comparison.

\begin{lemma}\label{lem:transfer}
Over an $S5$ frame, God-likeness is stable: if $\Gd^{\mathrm{A}}(x)$ at $w$ and
$r\,w\,v$, then $\Gd^{\mathrm{A}}(x)$ at $v$.
\end{lemma}

\begin{proof}
Let $\varphi$ be given. If $\Pos(\varphi)$ at $v$ then, by A4 at $v$ and
symmetry, $\Pos(\varphi)$ at $w$; so $\Box\varphi(x)$ at $w$; and transitivity
carries that to $\Box\varphi(x)$ at $v$. Conversely, if $\Box\varphi(x)$ at $v$
then symmetry and transitivity give $\Box\varphi(x)$ at $w$, hence
$\Pos(\varphi)$ at $w$, hence $\Pos(\varphi)$ at $v$ by A4.
\end{proof}

\noindent
God-likeness is then an $\ess^{\mathrm{A}}$-essence of any God-like $x$: the
forward direction of the biconditional follows from A4 and reflexivity, the
backward direction from Lemma~\ref{lem:transfer}. The rest runs as in
Theorem~\ref{thm:T3}, with $\Pos(\NE^{\mathrm{A}})$ supplying necessary
existence and reflexivity bringing $\Box\NE^{\mathrm{A}}(g)$ down to
$\NE^{\mathrm{A}}(g)$. Lemma~\ref{lem:andboxgod} obtains the conclusion of
Lemma~\ref{lem:transfer} from A3$^{*}$ alone, which is what makes the whole
detour avoidable.

\subsection{Fitting: extensions in place of intensions}

Fitting's move is different and, to our eye, the more illuminating of the two.
In G\"odel's and Scott's systems $\Pos$ is a predicate on \emph{intensional}
properties $I \to W \to \Prp$. Fitting makes it a predicate on
\emph{extensions} --- plain, world-independent sets $S : I \to \Prp$ --- and
re-reads the definitions accordingly:
\[
  \Gd^{\mathrm{F}}(x) \;:=\; \forall S : I \to \Prp,\ \Pos(S) \to S\,x,
\]
with essence and necessary existence likewise quantifying over extensions.

This is exactly what disarms the collapse. The proof of
Theorem~\ref{thm:collapse} took an arbitrary sentence $p$ and smuggled it into
the essence clause as the constant \emph{intensional} property
$\psi_p = \lambda y\,w.\,p(w)$, which tracks $p$ from world to world. An
extension has no world argument. There is no extension that varies with $p$,
the substitution has no counterpart, and the derivation stops.

The extensional reading also makes essences very cheap, and this time without
the reservations of \S\ref{sec:shortcut}.

\begin{lemma}[Singletons are essences]\label{lem:singleton}
For any $g$ and $w$, the extension $\{g\}$ is a Fitting-essence of $g$ at $w$,
with no axioms and over any frame.
\end{lemma}

\begin{proof}
The first conjunct is $g \in \{g\}$. For the second, let $T$ be an extension
with $g \in T$, let $r\,w\,v$, and let $z$ exist at $v$ with $z \in \{g\}$.
Then $z = g$, so $z \in T$.
\end{proof}

\noindent
Nothing here is contestable: extensions in Fitting's system are arbitrary sets
of individuals, and $\{g\}$ is one. Contrast Lemma~\ref{lem:haecceity}, which
needs a property mentioning a world precisely because Scott's essences range
over intensions.

\begin{theorem}[Theorem 3, de re]\label{thm:fitT3}
A1a, A2, T2 and A5 entail $\Box\exists x\,(E(x) \wedge \Gd^{\mathrm{F}}_w(x))$,
\emph{over an arbitrary frame} --- the argument runs in $K$. Neither A1b nor
A4 is used.
\end{theorem}

\begin{proof}
Theorem 1 --- which needs A1a and A2 --- applied to T2 gives a world accessible
from $w$ and an individual $g$ existing there with $\Gd^{\mathrm{F}}(g)$ as
evaluated at $w$. By A5, $g$ belongs to the necessary-existence extension at
$w$; applying that to $\{g\}$, which is an essence of $g$ by
Lemma~\ref{lem:singleton}, gives that $g$ itself exists at every world
accessible from $w$.
\end{proof}

\noindent
A4 was never used even by the argument through
Theorem~\ref{thm:T2}; A1b was, and Lemma~\ref{lem:singleton} removes the need
for it. The de dicto reading is a different matter.

\begin{theorem}[Theorem 3, de dicto]\label{thm:fitdd}
Adding A1b and A4, the same hypotheses entail
$\Box\exists x\,(E(x) \wedge \Gd^{\mathrm{F}}(x))$, with God-likeness
re-evaluated at each world. Also in $K$.
\end{theorem}

\begin{proof}
A4 carries the extension $\Gd^{\mathrm{F}}_w$ forward along $r$, so a God-like
individual at an accessible world is God-like at $w$; monotheism, which uses
A1b, identifies it with the witness supplied by Theorem~\ref{thm:fitT3}.
\end{proof}

\noindent
This gives a clean division of labour, and it is visible in the types of the
two theorems rather than only in their proofs: \textbf{A5 drives necessary
existence de re; A4 drives the world-stability of God-likeness; A1b enters only
through monotheism, which stability needs in order to identify witnesses.}
Together they yield necessary existence de dicto.

\noindent
The absence of any frame condition here is not an accident, and
Proposition~\ref{prop:nosymm} explains it. Symmetry was needed in Scott's
system to carry a conclusion from a merely possible world back to the actual
one. In Fitting's, the predicate $\Gd^{\mathrm{F}}(\cdot)$ evaluated at $w$ is
determined by the positivity facts \emph{at $w$}; a witness found at an
accessible world is a witness simpliciter, and there is nothing to carry back.
Freezing the predicate is what removes the need for the $B$ schema.

\subsection{H\'ajek: one axiom in place of two}\label{sec:hajek}

H\'ajek's route is different again. He merges A1 and A2 into a single weaker
axiom,
\[
  \text{H:A12}\qquad
  \bigl(\Pos(\varphi) \wedge \Box\forall x(\varphi(x) \to \psi(x))\bigr)
  \;\to\; \neg\,\Pos(\neg\psi),
\]
of which Anderson's A1a is the case $\psi := \varphi$, and which does
\emph{not} assert that $\psi$ is positive --- only that $\neg\psi$ is not. He
then rebuilds God-likeness on the derived notion
\[
  \Pos^{\#}(\varphi) \;:=\; \exists\psi\,\bigl(\Pos(\psi) \wedge \Box\forall y(\psi(y) \to \varphi(y))\bigr),
  \qquad
  \Gd^{\mathrm{H}}(x) \;:=\; \forall\varphi\,\bigl(\Box\varphi(x) \leftrightarrow \Pos^{\#}(\varphi)\bigr),
\]
keeping Anderson's essence and necessary existence unchanged. The axioms are
H:A12, A3, A4 and A5.

The striking feature is that the last two are not needed.

\begin{theorem}\label{thm:hajT3}
{\rm H:A12} and {\rm A3} alone, over a symmetric frame, entail
$\Box\exists x\,\Gd^{\mathrm{H}}(x)$. Neither {\rm A4} nor {\rm A5} is used.
\end{theorem}

\begin{proof}
Theorem 1 goes through from H:A12: if $\varphi$ were necessarily empty its
entailment hypothesis would hold vacuously for every $\psi$, and $\psi := \neg\varphi$
makes the conclusion say that $\varphi$ is not positive. Applied to A3 this
gives a world $v$ accessible from $w$ with some $g$ satisfying
$\Gd^{\mathrm{H}}(g)$ there.

Now the key step. $\Gd^{\mathrm{H}}$ is positive by A3, so it is necessarily
implied by a positive property --- itself --- and hence
$\Pos^{\#}(\Gd^{\mathrm{H}})$ holds outright. The biconditional in
$\Gd^{\mathrm{H}}(g)$ converts that directly into $\Box\Gd^{\mathrm{H}}(g)$:
the God-like individual is necessarily God-like. So
$\Box\exists x\,\Gd^{\mathrm{H}}(x)$ holds at $v$, symmetry brings
$\exists x\,\Gd^{\mathrm{H}}(x)$ back to $w$, and the same step applied at $w$
finishes.
\end{proof}

\noindent
Where Scott needs A5 and an essence argument to make God-like existence
necessary, H\'ajek's definition delivers it outright. Benzm\"uller, Weber and
Woltzenlogel Paleo report exactly this superfluousness, and add that A4 and A5
are nonetheless \emph{independent} of the rest --- superfluous without being
redundant~\cite{bww2017}. We have formalized the superfluousness; the
independence needs a further pair of models and is not done here.

This sits beside Corollary~\ref{cor:a5} in a way that invites a mistake worth
heading off. H\'ajek's A5 is stated with \emph{Anderson's} necessary existence,
and Proposition~\ref{prop:nevac} shows that $\NE^{\mathrm{A}}$ holds of every
individual at every world. One might conclude that A5 must be derivable here
too. It is not, and the reason is that A5 does not assert $\NE^{\mathrm{A}}$
but $\Pos(\NE^{\mathrm{A}})$. Anderson reaches the latter because A2 transfers
positivity along a necessary entailment whose consequent is true outright;
H\'ajek has no A2, and H:A12 concludes only that something \emph{fails} to be
positive, never that something is. A property being trivial does not make
calling it positive a triviality --- which is, in miniature, the reason
$\Pos$ was left undefined in the first place.

\subsection{All three avoid modal collapse}

\begin{proposition}\label{prop:variantsnocollapse}
Each of the three variants has a two-world model over the universal relation
--- an $S5$ frame --- satisfying all of its axioms, in which modal collapse
fails.
\end{proposition}

\begin{proof}
For Fitting, take one individual existing at both worlds and call a set
positive exactly when the individual belongs to it. For Anderson and for
H\'ajek, take one individual and call a property positive exactly when the
individual has it at \emph{every} world. In each case the axioms are routine to
verify, and $\lambda u.\,(u = w_0)$ holds at $w_0$ but not at the other world.
\end{proof}

\noindent
Each model also certifies the consistency of the variant it belongs to. The
mechanisms differ: Fitting's essence clause cannot express a contingent
sentence at all, whereas Anderson's and H\'ajek's God-likeness are stated
entirely in terms of $\Box$, so a contingent sentence never reaches their
biconditionals.

\begin{remark}
One small thing the formalization turned up. In Scott's system the step ``what
a God-like individual has is positive'' (Lemma~\ref{lem:posgod}) is classical,
because it argues by contradiction through A1b. Anderson has no A1b, and his
corresponding steps --- Lemma~\ref{lem:andboxgod}, Lemma~\ref{lem:transfer}, the
essence theorem, and Proposition~\ref{prop:nevac} with its
Corollary~\ref{cor:a5} --- are all \emph{constructive}; Lean reports them as
depending on no axioms at all. Fitting keeps A1b and with it the classical
step.

Classicality enters Anderson's system at exactly one place: Theorem~1, whose
proof is a reductio. Everything that routes through it inherits that, including
Theorem~\ref{thm:andT3} and Proposition~\ref{prop:a4}. It is worth recording
that Proposition~\ref{prop:mono} moved into this class when we removed its
reflexivity hypothesis: the earlier proof used reflexivity to produce an
accessible world and was constructive, whereas the present one obtains that
world from Theorem~1 and is not. Dropping a frame condition was not free; it
was paid for in logical strength elsewhere. Which of the two is preferable
depends on what one wants the certificate to certify, and we keep the weaker
hypothesis on the view that frame conditions are the scarcer resource here.
\end{remark}

\subsection{What is simplified here}\label{sec:variantscaveats}

Fitting's own setting is an intensional higher-order modal logic with actualist
quantifiers, and it supports distinctions this note does not draw. Two
departures should be recorded.

First, \emph{actualist quantifiers}. Our Fitting formalization carries an
explicit existence predicate and quantifies only over individuals existing at
the world in question. This is not optional: because a Fitting extension is
world-independent, with constant domains ``$S$ is empty'' would be a
world-independent statement, Theorem 1 would yield actual rather than merely
possible exemplification, and Theorem~\ref{thm:fitT3} would follow without ever
using A5 --- which is not Fitting's argument. Our Anderson formalization, by
contrast, uses constant domains, and there this costs nothing structural: his
God-likeness is intensional, so $\exists x\,\Gd^{\mathrm{A}}(x)$ still varies
from world to world and A5 still does real work.

Second, Fitting's logic has intension and extension operators which we have not
embedded, and which are needed to write his \emph{syntax}. As
Section~\ref{sec:derededicto} shows, the distinction those operators are for is
nonetheless available here, because in a shallow embedding it amounts to where
a world variable is bound.

\subsection{De re and de dicto}\label{sec:derededicto}

The predicate $\Gd^{\mathrm{F}}(z)$ carries two worlds' worth of information:
the individual $z$, and the world at which positivity is read off. So a claim
about God-like existence has two readings,
\[
  \underbrace{\forall v\,\bigl(r\,w\,v \to \exists z\ (\mathrm{E}z\,v \wedge \Gd^{\mathrm{F}}(z)\text{ at }w)\bigr)}_{\text{de re}}
  \qquad
  \underbrace{\forall v\,\bigl(r\,w\,v \to \exists z\ (\mathrm{E}z\,v \wedge \Gd^{\mathrm{F}}(z)\text{ at }v)\bigr)}_{\text{de dicto}}
\]
and no extra machinery is needed to state them: the difference is whether the
last world argument is $w$ or $v$. Theorem~\ref{thm:fitT3} is the de re
reading.

\begin{theorem}\label{thm:fitdd}
Fitting's axioms entail the de dicto reading as well, again over an arbitrary
frame.
\end{theorem}

\begin{proof}
It suffices that God-likeness is stable: if $\Gd^{\mathrm{F}}(x)$ at $w$ and
$r\,w\,v$, then $\Gd^{\mathrm{F}}(x)$ at $v$. By A4 the extension
$\gd_w := \{z : \Gd^{\mathrm{F}}(z)\text{ at }w\}$, positive at $w$ by T2,
is positive at $v$; so anything God-like at $v$ lies in $\gd_w$, that is, is
God-like at $w$. Theorem 1 supplies some $z$ God-like at $v$; that $z$ is then
God-like at $w$, and God-like individuals are unique at a world --- by the
argument of Proposition~\ref{prop:mono}, which works verbatim for Fitting's
system --- so $z = x$ and $x$ is God-like at $v$.
\end{proof}

\noindent
So with the full axiom set both readings hold. They do come apart, but earlier:
on the Part I axioms A1a, A1b, A2 and T2, before A4 is available. This is the
separation Benzm\"uller and Fuenmayor report, and it is worth being precise
about its location, since it is easy to mistake for a limitation on the
conclusion.

\begin{proposition}\label{prop:dd}
There is a model of the Part I axioms in which God-like existence is de re
possible but not de dicto possible.
\end{proposition}

\begin{proof}
Two worlds $\{0,1\}$ and two individuals $\{0,1\}$. Each world sees only the
other, $r\,w\,v \iff w \neq v$; the individual existing at a world is the other
one, $\mathrm{E}\,z\,v \iff z \neq v$; and a set is positive at $w$ exactly
when it contains $w$'s own individual.

Then $\Gd^{\mathrm{F}}(z)$ at $w$ holds iff $z = w$: positivity at $w$ is the
principal ultrafilter at $w$, so the only member of every positive set is $w$
itself. A1a, A1b and T2 are immediate. For A2, the world $\lnot w$ is
accessible from $w$ and $w$ exists there, so the entailment hypothesis
instantiates to $S\,w \to T\,w$.

Now $\Gd^{\mathrm{F}}$ is de re possible at $w$: the accessible world $\lnot w$
holds the individual $w$, which exists there and is God-like as read off at
$w$. But it is not de dicto possible, anywhere: the individual God-like at a
world is that world's own, and that is precisely the one that does not exist
there.
\end{proof}

\noindent
A4 must fail in this model, and does --- the singleton $\{0\}$ is positive at
world $0$ and not at world $1$. Given Theorem~\ref{thm:fitdd} it could not be
otherwise.

Two features of the separation are worth drawing out. It needs a
\emph{non-universal} frame: A4 propagates positivity along $r$, so when every
world sees every world it makes positivity world-independent, God-likeness
world-independent with it, and the two readings identical. And it needs
\emph{varying domains}: the room for them to come apart is an individual that
is God-like as read off at $w$ but does not exist at $v$. That is a second
reason, beyond the one in Section~\ref{sec:variantscaveats}, why the existence
predicate in this development is not optional.

\section{Frame conditions, tabulated}\label{sec:frames}

Keeping $r$ unconstrained in the semantics makes the following visible. Each
row lists the frame condition the result actually consumes, which in two cases
is weaker than the logic in which the result is usually stated.

\begin{center}
\begin{tabular}{@{}lll@{}}
\toprule
Result & Axioms used & Frame condition \\
\midrule
Theorem 1 (Lemma~\ref{lem:T1})        & A1a, A2                 & none \\
Empty Essence (Lemma~\ref{lem:empty}) & none                    & none \\
Inconsistency (Thm.~\ref{thm:main})   & A1a, A2, $\Pos(\NE)$    & none \quad($K$) \\
Theorem 2 (Thm.~\ref{thm:T2})         & A1b, A4                 & none \\
Seriality of $r$                      & A1a, A2, A3             & none \\
Monotheism                            & A1b                     & none \\
Theorem 3, Scott's route              & all of Scott's          & symmetry \quad($KB$) \\
Theorem 3 (Thm.~\ref{thm:T3})$^{\ast}$ & A1a, A2, A3, A5        & symmetry \quad($KB$) \\
Modal collapse (Thm.~\ref{thm:collapse})$^{\ast}$ & A1a, A2, A3, A5 & symmetry \quad($KB$) \\
\quad relative to T3 (Thm.~\ref{thm:mcrel}) & A1b, A4, serial, T3 & symmetry \quad($KB$) \\
$r$ is equality (Thm.~\ref{thm:frame})$^{\ast}$ & A1a, A2, A3, A5 & symmetry \quad($KB$) \\
One world (Thm.~\ref{thm:oneworld})$^{\ast}$ & A1a, A2, A3, A5   & universal \quad($S5U$) \\
\midrule
Anderson's Thm.~3 (Thm.~\ref{thm:andT3}) & A1a, A2, A3$^{*}$ \emph{only} & symmetry \quad($KB$) \\
\quad A5 for Anderson (Cor.~\ref{cor:a5}) & A2, A3$^{*}$        & none \quad($K$) \\
\quad A4 for Anderson (Prop.~\ref{prop:a4}) & A1a, A2, A3$^{*}$ & symm.\ + trans.\ \quad($K4B$) \\
Fitting's Thm.~3 de re (Thm.~\ref{thm:fitT3}) & A1a, A2, T2, A5 & none \quad($K$) \\
\quad de dicto (Thm.~\ref{thm:fitdd}) & \ $+$ A1b, A4            & none \quad($K$) \\
H\'ajek's Thm.~3 (Thm.~\ref{thm:hajT3})  & H:A12, A3 \emph{only} & symmetry \quad($KB$) \\
\bottomrule
\end{tabular}
\end{center}

\noindent
$^{\ast}$ These rows depend on Lemma~\ref{lem:haecceity}, and so on the
property domain containing world-indexed haecceities --- a question about the
semantics of the property quantifier, discussed there. Under full intensional
semantics they stand as they are; under a restricted property domain they may
not, and the unstarred route is the one the published analyses give.

Every entry is what our proof consumes. In an earlier version the Anderson row
carried a dagger, because our derivation then used $S5$ and the whole axiom set
while the published values were $KB$ with A4 and A5 redundant. The dagger is
gone: Theorem~\ref{thm:andT3} is now the $KB$ derivation, and the two
redundancy results are Corollary~\ref{cor:a5} and Proposition~\ref{prop:a4}.
The published results came first and were obtained by automated provers; ours
add no mathematics, only the guarantee that the stated hypotheses are the ones
actually used.

By Proposition~\ref{prop:nosymm} the symmetry entries for Scott are sharp: they
fail over a reflexive transitive frame, so no weakening in the $S4$ direction
is available there. The last three rows are the point of
Section~\ref{sec:variants}: Fitting pays no frame strength at all, H\'ajek pays
the same as Scott but with two of his four axioms unused, and Anderson pays the
same as Scott with two of his five axioms not merely unused but derivable. All
three avoid the collapse, at differing frame strengths. Frame strength is not
the variable that matters; the shape of the essence clause is.

The Anderson and H\'ajek rows are strikingly alike, and for the same structural
reason: both define God-likeness by a biconditional whose right-hand side is
already a $\Box$, so that the axiom asserting God-likeness positive converts in
one step into the necessity of God-likeness. Everything the essence apparatus
was built to deliver is then already in hand. That is why A4 and A5 --- the
rigidity of positivity and the positivity of necessary existence --- end up
carrying nothing in either system.

\noindent
Two comments. First, the refutation of G\"odel's axioms needs no frame
condition at all, so it is not an artefact of choosing $S5$; the axioms are
inconsistent in the weakest normal modal logic. Second, the positive argument
needs only symmetry, not the reflexivity and transitivity of $S5$ --- which
accords with the report in~\cite{bwp2016} that the argument goes through in
$KB$. Theorem~\ref{thm:main} is the one place where the accounting is subtle:
the refutation that Leo-II found automatically does use symmetry, and so lives
in $KB$; the human-constructed refutation of~\cite{bwp2016} formalized here
does not.

\section{The Lean 4 development}\label{sec:lean}

The development is 1461 lines across seven files, of which 378 are code and the
rest documentation, with \textbf{no dependency on Mathlib or any other library} --- nothing beyond Lean's core
logic is needed --- and builds in a couple of seconds. The shallow embedding is
direct:

\begin{quote}\small
\begin{verbatim}
abbrev Sentence (W : Type u) : Type u := W -> Prop
abbrev Property (I : Type v) (W : Type u) := I -> W -> Prop

def box (r : W -> W -> Prop) (p : Sentence W) : Sentence W :=
  fun w => forall v, r w v -> p v
\end{verbatim}
\end{quote}

\noindent
Each axiom is a named predicate (\verb|AxiomA1|, \verb|AxiomA1b|, \dots), so a
theorem's hypotheses state exactly which axioms it consumes; the two axiom sets
are bundled as \verb|Godel1970 r P| and \verb|Scott1987 r P|. The two headline
statements are

\begin{quote}\small
\begin{verbatim}
theorem inconsistent_of_A1_A2_A5 [Nonempty W]
    (hA1 : AxiomA1 P) (hA2 : AxiomA2 r P) (hA5 : AxiomA5 r P) : False

theorem modal_collapse (hs : Symm r) (h : Scott1987 r P)
    (p : Sentence W) (w : W) (hp : p w) : box r p w
\end{verbatim}
\end{quote}

\noindent
The first takes no hypothesis on \verb|r| at all, and the second takes exactly
one: that is how the middle column of Section~\ref{sec:frames} is recorded, and
it is checked rather than asserted. The correspondence with this note is

\begin{quote}\small
\begin{tabular}{@{}ll@{}}
\verb|possibly_exists|            & Lemma~\ref{lem:T1} \\
\verb|emptyProperty_ess|          & Lemma~\ref{lem:empty} \\
\verb|godel_1970_inconsistent|    & Theorem~\ref{thm:main} \\
\verb|empty_not_scottEss|         & Proposition~\ref{prop:scott} \\
\verb|positive_of_god|            & Lemma~\ref{lem:posgod} \\
\verb|scott_T2|                   & Theorem~\ref{thm:T2} \\
\verb|box_exists_god_of_god|      & Lemma~\ref{lem:boxgod} \\
\verb|exists_god|                 & Lemma~\ref{lem:existsgod} \\
\verb|scott_T3_viaT2|             & Theorem~\ref{thm:T3}, Scott's route \\
\verb|hae_scottEss|               & Lemma~\ref{lem:haecceity} \\
\verb|blind_of_god|               & Lemma~\ref{lem:blind} \\
\verb|scott_T3|                   & Theorem~\ref{thm:T3}, short route \\
\verb|serial_of_scott|            & seriality, over \verb|ScottCore| \\
\verb|modal_collapse|             & Theorem~\ref{thm:collapse} \\
\verb|modal_collapse_of_T3|       & Proposition~\ref{thm:mcrel} \\
\verb|s5_unique_world|            & Theorem~\ref{thm:oneworld} \\
\verb|accessibility_eq|           & Theorem~\ref{thm:frame} \\
\verb|box_iff|, \verb|dia_iff|    & Corollary~\ref{cor:inert} \\
\verb|toy_satisfies_A1_A2_A3_A4|, \verb|toy_not_A5| & Proposition~\ref{prop:toy} \\
\verb|toy_satisfies_scott|        & Proposition~\ref{prop:scottcons} \\
\verb|iso_satisfies_scott|        & Proposition~\ref{prop:iso} \\
\verb|ns_satisfies_scott|, \verb|ns_not_T3|, & \\
\quad \verb|ns_not_modal_collapse| & Proposition~\ref{prop:nosymm} \\
\verb|Anderson.box_god|           & Lemma~\ref{lem:andboxgod} \\
\verb|Anderson.monotheism|        & Proposition~\ref{prop:mono} \\
\verb|Anderson.T3|                & Theorem~\ref{thm:andT3} \\
\verb|Anderson.ne_universal|      & Proposition~\ref{prop:nevac} \\
\verb|Anderson.A5_of_A2_A3p|      & Corollary~\ref{cor:a5} \\
\verb|Anderson.A4_of_core|        & Proposition~\ref{prop:a4} \\
\verb|Anderson.god_transfers|     & Lemma~\ref{lem:transfer} \\
\verb|Fitting.singleton_ess|      & Lemma~\ref{lem:singleton} \\
\verb|Fitting.T3|                 & Theorem~\ref{thm:fitT3} \\
\verb|Fitting.T3_deDicto|         & Theorem~\ref{thm:fitdd} \\
\verb|Hajek.T3|                   & Theorem~\ref{thm:hajT3} \\
\verb|Anderson.noCollapse|, \verb|Fitting.noCollapse|, & \\
\quad \verb|Hajek.noCollapse|     & Proposition~\ref{prop:variantsnocollapse} \\
\end{tabular}
\end{quote}

Lean's \verb|#print axioms| reports what each theorem ultimately rests on. Every
classical result in the development reports exactly
\verb|[propext, Classical.choice, Quot.sound]|, and the following report
\emph{no axioms at all}:

\begin{quote}\small
\begin{verbatim}
emptyProperty_ess          does not depend on any axioms
empty_not_scottEss         does not depend on any axioms
toy_satisfies_A1_A2_A3_A4  does not depend on any axioms
toy_not_A5                 does not depend on any axioms
toy_satisfies_scott        does not depend on any axioms
\end{verbatim}
\end{quote}

Three things should be read off this. First, neither G\"odel's nor Scott's
axioms appear anywhere, because they are not axioms of the development: they
are hypotheses, discharged by whoever claims to have a model. Every theorem is
a conditional, which is what claims of this kind ought to be. Had the axioms
been declared with Lean's \verb|axiom| command, the development would have been
globally inconsistent and the certificate worthless.

Second, the only axioms that do appear are Lean's own three --- propositional
extensionality, choice, and soundness of quotients --- entering through the two
classical steps, Lemma~\ref{lem:T1} and Lemma~\ref{lem:posgod}.

Third, and most usefully, the consistency of Scott's system
(\verb|toy_satisfies_scott|) is certified constructively and axiom-free. It is
an explicit model, not the output of a model finder, and can be inspected.

\section{Scope of the result}\label{sec:caveats}

\paragraph{The comprehension assumption.}
This is the one place where the results are relative to a reading. The
quantifiers $\forall\varphi$ in the definitions of essence and of $\NE$ range
over \emph{all} properties $I \to W \to \Prp$, which is what makes the empty
property available for instantiation in Theorem~\ref{thm:main} --- and, on the
other side, what makes the constant property $\psi_p$ available in
Theorem~\ref{thm:collapse}. This is the standard second-order reading. A
defender could restrict the range of property quantification so as to exclude
necessarily uninstantiated properties, blocking Theorem~\ref{thm:main}; but
G\"odel states no such restriction, and any restriction strong enough to
exclude $e$ would need independent motivation, on pain of being introduced for
no reason except to save the argument. Scott's emendation blocks the derivation
without any such manoeuvre, which is a point in its favour --- though note that
a comprehension restriction would be no help against modal collapse, which
needs only a constant property.

\paragraph{What is not shown.}
Nothing here bears on whether God exists. Theorem~\ref{thm:main} says that one
particular list of premises is self-contradictory and therefore proves
everything, its intended conclusion included; an inconsistent argument is not a
disproof of its conclusion, it is the absence of an argument.
Theorem~\ref{thm:T3} says that another list of premises entails
$\Box\exists x\,\Gd(x)$; whether to accept those premises is a question about
the premises, and Theorem~\ref{thm:collapse} is the main reason most readers do
not.

\paragraph{Open.}
H\'ajek's emendations are untouched, and Fitting's intension and extension
operators are not embedded --- though Section~\ref{sec:derededicto} shows that
the distinction they exist to express is available without them. What a proper
embedding would add is the ability to write Fitting's own syntax, and with it
the remaining readings he distinguishes, rather than the two that matter for
the conclusion.

\section*{Availability}

The Lean sources, the build script and the axiom certificate are at
\url{https://github.com/alozanoroble/godel-ontological-lean}. The development
builds against Lean 4 \texttt{v4.34.0} with no external dependencies.

\begin{thebibliography}{99}

\bibitem{anderson}
C. A. Anderson.
\newblock Some emendations of G\"odel's ontological proof.
\newblock \emph{Faith and Philosophy} 7 (1990), no.~3, 291--303.

\bibitem{bwp2016}
C. Benzm\"uller and B. Woltzenlogel Paleo.
\newblock The inconsistency in G\"odel's ontological argument: a success story
  for AI in metaphysics.
\newblock In \emph{Proceedings of IJCAI-16}, pages 936--942, 2016.

\bibitem{bf2020}
C. Benzm\"uller and D. Fuenmayor.
\newblock Computer-supported analysis of positive properties, ultrafilters and
  modal collapse in variants of G\"odel's ontological argument.
\newblock \emph{Bulletin of the Section of Logic} 49 (2020), no.~2, 127--148.

\bibitem{bs2025}
C. Benzm\"uller and D. Scott.
\newblock Notes on G\"odel's and Scott's variants of the ontological argument.
\newblock \emph{Monatshefte f\"ur Mathematik} 208 (2025), 569--611.

\bibitem{bww2017}
C. Benzm\"uller, L. Weber and B. Woltzenlogel Paleo.
\newblock Computer-assisted analysis of the Anderson--H\'ajek ontological
  controversy.
\newblock \emph{Logica Universalis} 11 (2017), 139--151.

\bibitem{benz2026}
C. Benzm\"uller.
\newblock A comment on modal collapse and ultrafilters in G\"odel's ontological
  argument.
\newblock arXiv:2608.07578, 2026.

\bibitem{bk2026}
C. Benzm\"uller and D. Kirchner.
\newblock First-order modal logic in HOL: deep and shallow embeddings with
  automated faithfulness.
\newblock arXiv:2607.10880, 2026.

\bibitem{bkmso}
C. Benzm\"uller and D. Kirchner.
\newblock Monadic second-order logic in HOL: deep and shallow with automated
  faithfulness.
\newblock arXiv:2609.07345, 2026.

\bibitem{fitting}
M. Fitting.
\newblock \emph{Types, Tableaus, and G\"odel's God}.
\newblock Kluwer, 2002.

\bibitem{friedman}
H. M. Friedman.
\newblock A divine consistency proof for mathematics.
\newblock In \emph{Ontology of Divinity}. De Gruyter, 2024.

\bibitem{afp}
D. Fuenmayor and C. Benzm\"uller.
\newblock Types, tableaus and G\"odel's God in Isabelle/HOL.
\newblock \emph{Archive of Formal Proofs}, 2017.

\bibitem{goedel}
K. G\"odel.
\newblock Appendix A: notes in Kurt G\"odel's hand.
\newblock In J. H. Sobel, editor, \emph{Logic and Theism: Arguments for and
  Against Beliefs in God}, pages 144--145. Cambridge University Press, 2004.
\newblock Also in \emph{Collected Works, Volume III: Unpublished Essays and
  Lectures}, Oxford University Press, 1995.

\bibitem{hajek}
P. H\'ajek.
\newblock A new small emendation of G\"odel's ontological proof.
\newblock \emph{Studia Logica} 71 (2002), no.~2, 149--164.

\bibitem{hazen}
A. P. Hazen.
\newblock On G\"odel's ontological proof.
\newblock \emph{Australasian Journal of Philosophy} 76 (1998), no.~3, 361--377.

\bibitem{kovac}
S. Kova\v{c}.
\newblock Some weakened G\"odelian ontological systems.
\newblock \emph{Journal of Philosophical Logic} 32 (2003), no.~6, 565--588.

\bibitem{og2026}
P. Odifreddi and T. Gomes.
\newblock G\"odel's ontological proof and its consistency consequences.
\newblock \emph{Logica Universalis}, 2026.

\bibitem{scott}
D. S. Scott.
\newblock Appendix B: notes in Dana Scott's hand.
\newblock In J. H. Sobel, editor, \emph{Logic and Theism: Arguments for and
  Against Beliefs in God}, pages 145--146. Cambridge University Press, 2004.

\bibitem{sobel}
J. H. Sobel.
\newblock G\"odel's ontological proof.
\newblock In \emph{On Being and Saying: Essays for Richard Cartwright}, pages
  241--261. MIT Press, 1987.

\bibitem{wernhard}
C. Wernhard.
\newblock Applying second-order quantifier elimination in inspecting G\"odel's
  ontological proof.
\newblock arXiv:2110.11108, 2021.

\end{thebibliography}

\end{document}
